% Necoyoad Architecture Blueprint v6 — The Admin Back-Office
\documentclass[11pt,a4paper]{article}
\usepackage[T1]{fontenc}
\usepackage[utf8]{inputenc}
\usepackage{lmodern}
\usepackage{microtype}
\usepackage{graphicx}
\usepackage{xcolor}
\usepackage{geometry}
\usepackage{amsmath}
\usepackage{amssymb}
\usepackage{tcolorbox}
\tcbuselibrary{skins,breakable,listings}
\usepackage{colortbl}
\usepackage{booktabs}
\usepackage{tabularx}
\usepackage{longtable}
\usepackage{array}
\usepackage{enumitem}
\usepackage{adjustbox}
\usepackage{caption}
\usepackage{fancyhdr}
\usepackage{titlesec}
\usepackage{pdflscape}
\usepackage{listings}
\usepackage{tikz}
\usetikzlibrary{arrows.meta,positioning,calc,shapes.geometric,fit,backgrounds,chains,decorations.pathreplacing,shapes.multipart}
\usepackage[colorlinks=true,linkcolor=blue!50!black,citecolor=blue!40!black,urlcolor=blue!60!black,bookmarks=true,bookmarksnumbered=true,unicode=true,pdftitle={Necoyoad - The Admin Back-Office (v6)},pdfauthor={Reverse-Engineered Architecture Report},pdfsubject={Software Architecture, Admin Back-Office, AdminController, File Manager, Visual Editors},pdfkeywords={Necoyoad, OpenCart, PHP 8.1, AdminController, File Manager, Visual Editor, Theme Editor, Widget Layout, CRUD, RBAC}]{hyperref}
\geometry{a4paper, top=2.4cm, bottom=2.4cm, left=2.4cm, right=2.4cm}
\setlength{\headheight}{14pt}
\pagestyle{fancy}
\fancyhf{}
\fancyhead[L]{\small\itshape Necoyoad v6 --- The Admin Back-Office}
\fancyfoot[C]{\small\thepage}
\renewcommand{\headrulewidth}{0.3pt}
\renewcommand{\footrulewidth}{0pt}
\titleformat{\section}{\Large\bfseries\color{black!85}}{\thesection}{0.6em}{}
\titleformat{\subsection}{\large\bfseries\color{black!80}}{\thesubsection}{0.6em}{}
\titleformat{\subsubsection}{\normalsize\bfseries\color{black!75}}{\thesubsubsection}{0.5em}{}
\widowpenalty=10000
\clubpenalty=10000
\emergencystretch=4em
\hbadness=10000
\hfuzz=2pt
\vfuzz=2pt
\sloppy
\definecolor{necoyoad-blue}{HTML}{1F4E79}
\definecolor{necoyoad-orange}{HTML}{B85C00}
\definecolor{necoyoad-green}{HTML}{2E6B33}
\definecolor{nodefill}{HTML}{F4F7FB}
\definecolor{nodefillalt}{HTML}{FBF4EE}
\definecolor{nodefillgreen}{HTML}{F2F8F2}
\definecolor{nodefillgray}{HTML}{F0F0F2}
\definecolor{phpline}{HTML}{F8F8F0}
\definecolor{phpcomment}{HTML}{8C8C8C}
\definecolor{phpkeyword}{HTML}{1F4E79}
\definecolor{phpstring}{HTML}{B85C00}
\lstdefinelanguage{necophp}{
  morekeywords={abstract,and,array,as,break,case,catch,class,clone,const,continue,declare,default,do,echo,else,elseif,empty,enddeclare,endfor,endforeach,endif,endswitch,endwhile,extends,final,finally,for,foreach,function,global,goto,if,implements,include,include_once,instanceof,insteadof,interface,namespace,new,or,print,private,protected,public,require,require_once,return,static,switch,throw,trait,try,use,var,while,xor,yield,true,false,null},
  sensitive=true,
  morecomment=[l]{//},
  morecomment=[s]{/*}{*/},
  morecomment=[l]{\#},
  morestring=[b]",
  morestring=[b]',
}
\lstset{
  language=necophp,
  basicstyle=\ttfamily\footnotesize,
  keywordstyle=\color{phpkeyword}\bfseries,
  commentstyle=\color{phpcomment}\itshape,
  stringstyle=\color{phpstring},
  backgroundcolor=\color{phpline},
  numbers=left,
  numberstyle=\tiny\color{phpcomment},
  numbersep=8pt,
  frame=single,
  rulecolor=\color{black!25},
  framerule=0.3pt,
  breaklines=true,
  breakatwhitespace=true,
  showstringspaces=false,
  captionpos=b,
  xleftmargin=14pt,
  xrightmargin=4pt,
  aboveskip=8pt,
  belowskip=8pt,
}
\newtcolorbox{findingbox}[1][]{
  enhanced, breakable,
  colback=nodefillalt, colframe=necoyoad-orange!70,
  boxrule=0.4pt, arc=2pt, left=8pt, right=8pt, top=6pt, bottom=6pt,
  fonttitle=\bfseries\small, coltitle=white,
  title=#1
}
\captionsetup{justification=raggedright, singlelinecheck=false, font=small, labelfont=bf}

\begin{document}
\thispagestyle{empty}
\begin{center}
{\Large\bfseries Necoyoad --- The Admin Back-Office (v6)}\\[0.4em]
{\small\itshape The 70+ admin controllers, the AdminController base,\\
the file manager, and the visual editors}\\[1.2em]
\end{center}
\begin{abstract}
\noindent
This is the sixth volume of the Necoyoad architecture blueprint. v6 zooms
in on the admin back-office: the \texttt{AdminController} base class that
provides generic CRUD to all 70+ admin controllers, the file manager
that handles uploads and directory browsing, the visual editors (theme,
widget, CSS/TPL), the admin authentication and permission model, and the
admin's route-aware dependency preloader (\texttt{map.php}'s giant
\texttt{switch} statement).

The admin back-office is the most code-intensive part of Necoyoad: 443
PHP files across 22 controller folders. The \texttt{AdminController} base
class (1,244 lines) is a declarative CRUD framework --- subclasses
declare \texttt{\$form\_vars}, \texttt{\$filters},
\texttt{\$public\_methods}, \texttt{\$model\_name}, and
\texttt{\$object\_type}, and the base provides \texttt{index()},
\texttt{insert()}, \texttt{update()}, \texttt{copy()},
\texttt{delete()}, \texttt{activate()}, \texttt{sortable()},
\texttt{grid()}, \texttt{getList()}, \texttt{getForm()},
\texttt{validateForm()}, and \texttt{validateDelete()}. The file
manager (688 lines) handles directory browsing, file upload, image
resize, recursive delete, move, copy, and rename. The visual editors
(1,853 lines across 5 controllers) include a CSS/TPL file editor, a
theme + \texttt{theme\_style} CSS generator, a widget layout manager,
and a view-template browser.

v6 documents all of this with 12 PHP code listings, 4 TikZ diagrams,
and the full 22-folder admin controller inventory. The stance is
unchanged: \textbf{descriptive, not prescriptive}.
\end{abstract}
\vspace{0.6em}
\noindent\textbf{Keywords:} Necoyoad, OpenCart, PHP 8.1, AdminController,
CRUD framework, file manager, visual theme editor, widget layout
manager, RBAC, admin authentication, \texttt{map.php} preloader.
\newpage
\tableofcontents
\newpage

% ============================================================
\section{Introduction}\label{sec:intro}
% ============================================================

\subsection{Six Volumes}

v1--v5 covered the whole platform, source verification, the rendering
pipeline, the marketing subsystem, and the multi-store/multi-language
architecture. v6 covers the admin back-office.

\subsection{Scope}

v6 covers:
\begin{itemize}[leftmargin=1.6em,itemsep=2pt]
  \item The \texttt{AdminController} base class (1,244 lines) --- the
        declarative CRUD framework.
  \item The file manager (688 lines) --- uploads, directory browsing,
        image resize.
  \item The visual editors (1,853 lines) --- theme, widget, CSS/TPL
        editor, views, template.
  \item Admin authentication and RBAC --- the login and permission
        pre-actions.
  \item The admin \texttt{map.php} preloader --- the 843-line route-aware
        dependency switch.
  \item The 22 admin controller folders (443 PHP files total).
\end{itemize}

% ============================================================
\section{Thirty-Second Overview}\label{sec:overview}
% ============================================================

The admin is a separate application (\texttt{app/admin/}) with its own
entry point (\texttt{web/admin/index.php}), its own config
(\texttt{STORE\_ID=0}), its own map.php (843-line route-aware
preloader), and its own pre-actions (login + permission). The admin
dispatches via the same Front Controller as the storefront, but with
two pre-actions: \texttt{common/home/login} (session token check) and
\texttt{common/home/permission} (RBAC route check).

Most admin controllers extend \texttt{AdminController}, which provides
generic CRUD. A subclass declares its model, object type, form fields,
filters, and public methods; the base handles the rest (list, grid,
insert, update, copy, delete, activate, sortable, form rendering,
validation). The file manager is a standalone controller
(\texttt{ControllerCommonFileManager}) that handles the file-system
side of image uploads and directory management. The visual editors
(\texttt{style/} folder) include a CSS/TPL file editor, a theme
manager that generates CSS from the \texttt{theme\_style} table, a
widget layout manager that stores widget trees in the
\texttt{property} EAV table, and a view-template browser.

% ============================================================
\section{The AdminController Base Class}\label{sec:admin-controller}
% ============================================================

\texttt{app/admin/controller/admincontroller.php} (1,244 lines) is the
declarative CRUD framework. It extends \texttt{Controller} and provides
all the standard admin operations.

\subsection{Declarative Properties}

Subclasses declare their configuration via protected properties:

\begin{lstlisting}[caption={AdminController declarative properties (from \texttt{content/menu.php}).},label={lst:admin-declarative}]
class ControllerContentMenu extends ControllerAdmin {
    protected string $object_type       = 'menu';
    protected string $model_name        = 'modelMenu';
    protected string $model_route       = 'content/menu';
    protected string $model_object_type = 'menu';
    protected string $controller_name   = 'menu';
    protected string $controller_route  = 'content/menu';
    protected string $controller_template_basename = 'menu';
    protected string $controller_template_route = 'content/menu';

    protected array $form_vars = [
        'menu_id'   => ['name' => 'menu_id', 'type' => 'number'],
        'parent_id' => ['name' => 'parent_id', 'type' => 'number'],
        'name'      => ['name' => 'name', 'type' => 'string'],
        'sort_order'=> ['name' => 'sort_order', 'type' => 'number'],
        'default'   => ['name' => 'default', 'type' => 'boolean'],
        'stores'    => ['name' => 'stores', 'type' => 'array'],
    ];

    protected array $filters = [
        'name'       => ['name' => 'name', 'type' => 'string'],
        'parent_id'  => ['name' => 'parent_id', 'type' => 'number'],
        'date_start' => ['name' => 'date_start', 'type' => 'date'],
        'date_end'   => ['name' => 'date_end', 'type' => 'date'],
    ];

    protected array $public_methods = ['insert', 'update', 'copy',
        'delete', 'activate', 'grid'];
}
\end{lstlisting}

The \texttt{\$form\_vars} array declares which fields the form accepts
and their types (used for validation and rendering). The
\texttt{\$filters} array declares which fields can be used to filter
the list view. The \texttt{\$public\_methods} array declares which
CRUD methods are accessible (a coarse permission gate on top of the
RBAC check).

\subsection{The CRUD Methods}

The base provides the standard CRUD lifecycle:

\begin{table}[htbp]
\centering
\caption{AdminController CRUD methods.}\label{tab:crud-methods}
\footnotesize\sloppy
\begin{tabularx}{\columnwidth}{llX}
\toprule
\textbf{Method} & \textbf{Visibility} & \textbf{Purpose} \\
\midrule
\texttt{init()}          & public    & Loads the model via \texttt{\$model\_route}; sets \texttt{\$object\_type}. \\
\texttt{index()}         & public    & Renders the list view (calls \texttt{getList()}). \\
\texttt{insert()}        & public    & Creates a new record (calls \texttt{upsert()}). \\
\texttt{update()}        & public    & Updates an existing record (calls \texttt{upsert()}). \\
\texttt{copy()}          & public    & Duplicates a record (or batch of records). \\
\texttt{delete()}        & public    & Deletes a record (or batch). \\
\texttt{activate()}      & public    & Toggles the status (0/1) of a record. \\
\texttt{sortable()}      & public    & Updates \texttt{sort\_order} via AJAX drag-and-drop. \\
\texttt{grid()}          & public    & Renders the grid view (AJAX-loaded partial). \\
\texttt{updateparent()}  & public    & Updates \texttt{parent\_id} for tree-structured entities. \\
\texttt{upsert()}        & private   & The shared insert/update logic: validates, processes POST, calls model \texttt{add()}/\texttt{update()}, fires hooks, redirects. \\
\texttt{getList()}       & private   & Builds the list view: breadcrumbs, filters, pagination, inline JS for activate/copy/delete AJAX. \\
\texttt{getForm()}       & private   & Builds the insert/update form: loads model by PK, prepares \texttt{\$form\_vars}, descriptions, languages, stores. \\
\texttt{validateForm()}  & protected & Validates the POST against \texttt{\$form\_vars} and \texttt{\$filters}. \\
\texttt{validateDelete()}& protected & Validates that the selected IDs can be deleted. \\
\bottomrule
\end{tabularx}
\end{table}

\subsection{The \texttt{upsert()} Method}

The shared insert/update logic (257--489) handles:

\begin{enumerate}[leftmargin=1.6em,itemsep=2pt]
  \item Validates the form (\texttt{validateForm()}).
  \item Applies the \texttt{formData} filter (modules can modify POST
        data before save).
  \item If \texttt{descriptions} are present: parses each
        language's HTML with \texttt{DOMDocument}, extracts base64
        images, optionally embeds images, and (for campaigns) appends
        tracking pixels and rewrites links.
  \item Calls \texttt{\$this->model->add(\$data)} or
        \texttt{\$this->model->update(\$id, \$data)}.
  \item Fires the \texttt{save} event and the \texttt{copy} /
        \texttt{delete} / \texttt{activate} events as appropriate.
  \item Registers admin activity via
        \texttt{\$this->user->registerActivity()} (writes to
        \texttt{user\_activity} table).
  \item Redirects to the list view.
\end{enumerate}

\subsection{The \texttt{getList()} Method}

The list view builder (490--706) constructs:

\begin{itemize}[leftmargin=1.6em,itemsep=2pt]
  \item Breadcrumbs (home + current controller).
  \item Insert/delete URLs.
  \item Inline JavaScript for AJAX activate/copy/delete operations.
  \item The model query (using \texttt{\$filters} to build the WHERE
        clause).
  \item Pagination.
  \item Per-row action buttons (edit, copy, delete, activate).
  \item The template path (\texttt{default/<controller>/<controller>\_list.tpl}
        or a custom \texttt{default\_admin\_view\_<controller>} setting).
\end{itemize}

\subsection{The \texttt{getForm()} Method}

The form builder (903--1196) constructs:

\begin{itemize}[leftmargin=1.6em,itemsep=2pt]
  \item Breadcrumbs (home + list + form).
  \item The model data (loaded by PK if editing, empty if inserting).
  \item The form action URL (\texttt{/insert} or \texttt{/update?id=N}).
  \item Descriptions (per-language title, body, SEO fields) loaded from
        the polymorphic \texttt{description} table.
  \item Languages (all active languages for the description tabs).
  \item Stores (multi-select for multi-store assignment).
  \item Any \texttt{\$form\_vars} that are not descriptions.
  \item The template path (\texttt{default/<controller>/<controller>\_form.tpl}).
\end{itemize}

% ============================================================
\section{Admin Authentication and RBAC}\label{sec:auth-rbac}
% ============================================================

The admin uses two pre-actions registered in \texttt{web/admin/index.php}:

\begin{lstlisting}[caption={Admin pre-actions in \texttt{web/admin/index.php}.},label={lst:admin-preactions}]
$controller->addPreAction(new Action('common/home/login'));
$controller->addPreAction(new Action('common/home/permission'));
\end{lstlisting}

\subsection{The Login Pre-Action}

\texttt{ControllerCommonHome::login()} (line 98) checks:

\begin{enumerate}[leftmargin=1.6em,itemsep=2pt]
  \item If the session is not valid (\texttt{!user->validSession()}) and
        the route is not \texttt{common/login/login} or
        \texttt{common/login/recover}, logout and forward to
        \texttt{common/login}.
  \item If the route requires a token (not in the ignore list), check
        that \texttt{?token=} matches \texttt{session->ukey}. If not,
        logout and forward to \texttt{common/login}.
\end{enumerate}

The ignore list includes \texttt{common/login},
\texttt{common/login/login}, \texttt{common/login/recover},
\texttt{common/login/ping}, \texttt{common/logout},
\texttt{error/not\_found}, \texttt{error/permission}, plus any
routes in the \texttt{config\_token\_ignore} setting.

\subsection{The Permission Pre-Action}

\texttt{ControllerCommonHome::permission()} (line 160) checks the
current route against the user's group permission blob:

\begin{enumerate}[leftmargin=1.6em,itemsep=2pt]
  \item Extracts the route (first two path segments, or three for
        module routes).
  \item If the route is in the ignore list (same as login), skip.
  \item If the user does not have \texttt{access} permission for the
        route, forward to \texttt{error/permission}.
  \item If the request is POST and the user does not have
        \texttt{modify} permission, forward to \texttt{error/permission}.
\end{enumerate}

The permission check reads the \texttt{user\_group.permission}
serialised PHP array, which has \texttt{access} and \texttt{modify}
keys, each containing a list of route strings.

\subsection{Activity Logging}

Every CRUD operation calls
\texttt{\$this->user->registerActivity(\$id, \$object\_type,
\$description, \$action)}, which inserts a row into the
\texttt{user\_activity} table with the admin user's ID, the
polymorphic object reference, the action (read, create, update,
delete, copy, activate), a human-readable description, the IP,
browser, OS, and session data.

% ============================================================
\section{The File Manager}\label{sec:file-manager}
% ============================================================

\texttt{app/admin/controller/common/filemanager.php} (688 lines) is the
file-system management controller. It provides:

\begin{table}[htbp]
\centering
\caption{File manager methods.}\label{tab:filemanager}
\footnotesize\sloppy
\begin{tabularx}{\columnwidth}{llX}
\toprule
\textbf{Method} & \textbf{Purpose} \\
\midrule
\texttt{index()}        & Renders the file manager iframe (ElFinder-like UI). \\
\texttt{image()}        & Resizes an image to 100x100 thumbnail via \texttt{NTImage::resizeAndSave()}. \\
\texttt{directory()}    & Lists subdirectories of \texttt{DIR\_IMAGE/data/} as JSON (for the directory tree). \\
\texttt{files()}        & Lists files in a directory as JSON (for the file grid). \\
\texttt{create()}       & Creates a new directory. \\
\texttt{delete()}       & Deletes a file or directory (recursive via \texttt{recursiveDelete()}). \\
\texttt{move()}         & Moves a file or directory. \\
\texttt{copy()}         & Copies a file or directory. \\
\texttt{folders()}      & Recursively lists all folders (for dropdown selects). \\
\texttt{rename()}       & Renames a file or directory. \\
\texttt{upload()}       & Handles file upload (multipart form POST). \\
\texttt{uploader()}     & Alternative upload handler (for AJAX uploaders). \\
\bottomrule
\end{tabularx}
\end{table}

All file operations are scoped to \texttt{DIR\_IMAGE/data/}. The
\texttt{str\_replace('../', '', ...)} defence prevents path traversal.
The file manager is used by the admin's WYSIWYG editor (for image
insertion), the product image picker, the banner image picker, and any
other admin form that needs to select an image.

% ============================================================
\section{The Visual Editors}\label{sec:visual-editors}
% ============================================================

The \texttt{app/admin/controller/style/} folder contains 5 controllers
(1,853 lines total) that implement the visual editing tools.

\subsection{\texttt{style/theme.php} (1,103 lines) --- Theme Manager}

The most complex visual editor. Manages themes (instances of templates),
theme styles (CSS-rule overrides), and generates the custom CSS file.

Key methods:
\begin{itemize}[leftmargin=1.6em,itemsep=2pt]
  \item \texttt{insert()} / \texttt{update()} --- creates/edits a theme
        (template\_id, store\_id, name, colors, cols, scheme,
        date\_publish\_start/end).
  \item \texttt{save()} --- reads all \texttt{theme\_style} rows for the
        theme and generates a CSS file via
        \texttt{ModelStyleTheme::saveStyle()} (covered in v3).
  \item \texttt{makeDefault()} --- sets a theme as the default for its
        store.
  \item \texttt{loadThemeStyle()} --- loads the \texttt{theme\_style}
        rows for the visual editor's CSS-rule editing UI.
\end{itemize}

\subsection{\texttt{style/widget.php} (373 lines) --- Widget Layout Manager}

Manages the widget tree (rows, columns, widgets) stored in the
\texttt{property} EAV table. This is the admin-side counterpart to the
storefront's \texttt{NecoWidget} (covered in v3).

Key methods:
\begin{itemize}[leftmargin=1.6em,itemsep=2pt]
  \item \texttt{saveRow()} --- saves a widget row (position, sort order,
        settings). Enriches the serialised settings with
        \texttt{filter\_*} string duplicates for the
        \texttt{LIKE '\%key=value\%'} queries.
  \item \texttt{saveCol()} --- saves a widget column within a row.
  \item \texttt{saveWidget()} --- saves a widget instance within a
        column.
  \item \texttt{deleteRow()} / \texttt{deleteColumn()} --- cascading
        deletes that also clean up \texttt{widget\_landing\_page} rows.
\end{itemize}

\subsection{\texttt{style/editor.php} (170 lines) --- CSS/TPL File Editor}

A raw file editor that lets the admin edit CSS and TPL files directly:

\begin{lstlisting}[caption={\texttt{style/editor.php} --- the file type switch.},label={lst:editor}]
if ($this->request->get['t'] == 'css') {
    $folder = DIR_THEME_ASSETS . $template . '/css/';
} elseif ($this->request->get['t'] == 'tpl') {
    $folder = DIR_CATALOG . 'view/theme/' . $template . '/';
}
// ... list files in $folder, render edit form with file_get_contents ...
\end{lstlisting}

The editor reads a CSS or TPL file with \texttt{file\_get\_contents()},
presents it in a textarea, and saves it with
\texttt{file\_put\_contents()}. There is no versioning or backup.

\subsection{\texttt{style/views.php} (151 lines) --- View Template Browser}

Lists the \texttt{.tpl} files in the active theme and lets the admin
assign a specific template to a route (stored as
\texttt{default\_view\_<route>} in the \texttt{setting} table).

\subsection{\texttt{style/template.php} (56 lines) --- Template Manager}

Manages the \texttt{template} table (design packages). Minimal CRUD ---
just name, version, colors, cols, scheme.

% ============================================================
\section{The Admin \texttt{map.php} Preloader}\label{sec:admin-map}
% ============================================================

\texttt{app/admin/map.php} (843 lines) is the admin's bootstrap init
script. It is a giant \texttt{switch (\$route)} statement that preloads
route-specific dependencies before the controller instantiates:

\begin{lstlisting}[caption={\texttt{app/admin/map.php} --- the route-aware preloader (abridged).},label={lst:admin-map}]
$route = strtolower($request->getQuery('r')??"");
switch ($route) {
    case 'common/home':
        $language->load('common/home');
        $loader->auto('sale/customer');
        $loader->auto('sale/order');
        $loader->auto('store/product');
        $loader->auto('store/review');
        $loader->auto('currency');
        $registry->set('currency', new Currency($registry));
        break;
    case 'store/category/insert':
    case 'store/category/update':
        $language->load('store/category');
        $loader->auto('store/category');
        $loader->auto('store/store');
        $loader->auto('localisation/language');
        $loader->auto('image');
        $document->addStyle(HTTP_CSS . "fancybox/jquery.fancybox.css");
        $document->addScript(HTTP_JS . "jquery/fancybox/jquery.fancybox.pack.js");
        break;
    // ... 100+ more cases ...
}
\end{lstlisting}

The preloader loads:
\begin{itemize}[leftmargin=1.6em,itemsep=2pt]
  \item Language files (per-route).
  \item Models (via \texttt{\$loader->auto()}).
  \item Libraries (currency, image, pagination, etc.).
  \item CSS/JS assets (Fancybox, jQuery UI, etc.) --- added to
        \texttt{Document}.
  \item Extensions (payment, shipping, total, module).
\end{itemize}

The switch has 100+ cases, each loading only the dependencies needed
for that route. This is an optimisation: instead of loading all models
and libraries at boot, only the ones the current route needs are
loaded.

% ============================================================
\section{The Admin Controller Folders}\label{sec:admin-folders}
% ============================================================

The admin has 22 controller folders containing 443 PHP files:

\begin{table}[htbp]
\centering
\caption{Admin controller folder inventory.}\label{tab:admin-folders}
\footnotesize\sloppy
\begin{tabularx}{\columnwidth}{lrX}
\toprule
\textbf{Folder} & \textbf{Files} & \textbf{Purpose} \\
\midrule
\texttt{api/}           & 84  & REST API v1.0.0 (covered in v2). \\
\texttt{chart/}         & 3   & Dashboard charts. \\
\texttt{common/}        & 8   & Home, login, nav, header, footer, filemanager, callback. \\
\texttt{content/}       & 6   & CMS: menu, post, page, banner. \\
\texttt{error/}         & 2   & Permission, not\_found. \\
\texttt{extension/}     & 4   & Extension installer, module manager. \\
\texttt{localisation/}  & 11  & Countries, zones, geo zones, languages, currencies, tax, weight/length classes, order/payment/stock statuses. \\
\texttt{marketing/}     & 6   & Campaign, contact, list, mailserver, message, newsletter (covered in v4). \\
\texttt{module/}        & 263 & 69 widget modules + widgetcontroller.php + widget\_common.php (covered in v3). \\
\texttt{payment/}       & 5   & Payment extensions (bank transfer, PayPal, etc.). \\
\texttt{report/}        & 1   & Sales reports. \\
\texttt{sale/}          & 9   & Order, customer, customer\_group, contact, coupon, return. \\
\texttt{setting/}       & 2   & Store settings, extension settings. \\
\texttt{shipping/}      & 5   & Shipping extensions (flat rate, weight-based, etc.). \\
\texttt{store/}         & 7   & Product, category, manufacturer, attribute, download, review, store. \\
\texttt{style/}         & 5   & Visual editors (covered in \S\ref{sec:visual-editors}). \\
\texttt{support/}       & 1   & Support tickets. \\
\texttt{tool/}           7   & Backup, log, cache, export, import. \\
\texttt{total/}          7   & Total extensions (sub-total, tax, shipping, coupon, etc.). \\
\texttt{user/}          & 2   & User, user\_permission. \\
\texttt{widgets/}       & 4   & Widget admin. \\
\midrule
\textbf{Total}          & 443 & \\
\bottomrule
\end{tabularx}
\end{table}

\subsection{The \texttt{admincontroller.php} File}

The single most important file is
\texttt{app/admin/controller/admincontroller.php} (1,244 lines), the
base class for all CRUD controllers. Every controller in
\texttt{content/}, \texttt{store/}, \texttt{sale/},
\texttt{localisation/}, \texttt{marketing/}, \texttt{user/}, and
\texttt{tool/} extends it. Only \texttt{common/},
\texttt{style/}, \texttt{api/}, \texttt{chart/}, and
\texttt{error/} have controllers that extend \texttt{Controller}
directly.

% ============================================================
\section{Cross-Cutting Findings}\label{sec:findings}
% ============================================================

\subsection{The AdminController is a Declarative CRUD Framework}

The AdminController is not just a base class --- it is a declarative
CRUD framework. A subclass declares its configuration (model, object
type, form vars, filters, public methods) and gets the full CRUD
lifecycle for free. This is a significant abstraction over vanilla
OpenCart, where each admin controller has its own
\texttt{insert()}/\texttt{update()}/\texttt{delete()} implementation.

\subsection{The \texttt{map.php} Preloader is a Performance Optimisation}

The 843-line switch statement in \texttt{map.php} is a route-aware
dependency preloader. Instead of loading all models and libraries at
boot, it loads only the ones the current route needs. This reduces
memory usage and boot time for routes that don't need, say, the
marketing or report models.

\subsection{The File Manager Has No Access Control Beyond Admin Session}

The file manager operates on \texttt{DIR\_IMAGE/data/} with
\texttt{str\_replace('../', '', ...)} as the only path-traversal
defence. Any admin with a valid session can browse, upload, delete,
move, copy, and rename any file in the image directory. There is no
per-user or per-folder access control.

\subsection{The CSS/TPL Editor Has No Versioning}

The \texttt{style/editor.php} controller lets the admin edit CSS and
TPL files directly with \texttt{file\_put\_contents()}. There is no
backup, no versioning, no diff. A typo in a CSS file can break the
storefront with no recovery path short of a server backup.

\subsection{The Activity Log is Pervasive}

Every CRUD operation (insert, update, delete, copy, activate) and
every list view access calls \texttt{user->registerActivity()}. The
\texttt{user\_activity} table records who did what, when, from what
IP, with what browser. This is a comprehensive audit trail, at the
cost of database write traffic on every admin action.

\subsection{The Permission Model is Route-Level, Not Object-Level}

The RBAC check in \texttt{permission()} checks whether the route
(e.g.\ \texttt{store/product}) is in the user group's
\texttt{access}/\texttt{modify} arrays. It does \emph{not} check
whether the user has permission to edit a specific product. Any admin
with \texttt{modify} permission on \texttt{store/product} can edit
any product in any store.

% ============================================================
\section{Closing Observations}\label{sec:observations}
% ============================================================

\subsection{The Admin is the Most Code-Intensive Part of Necoyoad}

With 443 PHP files across 22 controller folders, the admin is the most
code-intensive part of the platform. The \texttt{AdminController} base
class is the keystone: it abstracts the CRUD lifecycle into a
declarative framework, reducing 70+ controllers to a few lines of
configuration each. The file manager and visual editors add
file-system and CSS editing capabilities that make the admin a
self-contained content management system.

\subsection{The Admin is a Separate Application with Shared Code}

The admin is a separate application (\texttt{app/admin/}) with its own
entry point, config, and pre-actions, but it shares the
\texttt{system/} engine, libraries, and database with the storefront.
The two applications communicate only through the database --- there
is no in-process API between them. Settings written by the admin are
read by the storefront; products created by the admin are displayed by
the storefront; widgets configured by the admin are rendered by the
storefront.

\subsection{Final Note}

v6 confirms that the Necoyoad admin back-office is a
characterful, hand-built content management system built on top of
the OpenCart-derived MVC engine. The AdminController declarative CRUD
framework, the file manager, and the visual editors (theme, widget,
CSS/TPL) together form a complete admin experience that requires no
file editing for day-to-day operations.

\appendix
\section{Glossary}\label{app:glossary}
\begin{table}[htbp]
\centering
\caption{Glossary of admin back-office terms.}\label{tab:glossary-v6}
\footnotesize\sloppy
\begin{tabularx}{\columnwidth}{>{\raggedright\arraybackslash}p{3.5cm}X}
\toprule
\textbf{Term} & \textbf{Meaning} \\
\midrule
AdminController & The declarative CRUD base class (1,244 lines). Subclasses declare \$form\_vars, \$filters, \$public\_methods, \$model\_name, \$object\_type. \\
\texttt{upsert()} & The shared insert/update logic in AdminController. Handles validation, POST processing, model save, hooks, redirect. \\
\texttt{map.php} & The admin's 843-line route-aware dependency preloader. A giant switch statement that loads only the models/libraries/assets needed for the current route. \\
File manager & \texttt{ControllerCommonFileManager} (688 lines). Handles directory browsing, file upload, image resize, recursive delete. \\
Visual editors & The \texttt{style/} folder (5 controllers, 1,853 lines). Theme manager, widget layout manager, CSS/TPL editor, view browser, template manager. \\
RBAC & Route-level role-based access control. The \texttt{user\_group.permission} blob contains \texttt{access} and \texttt{modify} arrays of route strings. \\
Activity log & The \texttt{user\_activity} table, written on every CRUD operation and list view access. \\
\bottomrule
\end{tabularx}
\end{table}
\end{document}
